% SGA 3，第 VII 篇_A，中文重构——第 1 轮。
% 范围：权威底本本地第 15--21 页，第 3 节全文，止于本地第 22 页第 4 节标题之前。
% 合并阅读本页码： 468--474. 再版本印刷页码： 457--463.
% 本段原版页边定位： 426--433.
% 权威底本：Polo--Gille Exp7A-13oct24.pdf,
% SHA-256 D46DB8C79D0385643B6EED4B96900613248375F9B73B4202478A4F57FD3E8522.
% 参照了 Jacob Reinhold 的 CC BY 4.0 Markdown 作为比对谱系。

\section*{3. 余代数与 Cartier 对偶}
\phantomsection
\label{sga3:VIIA:3}
\addcontentsline{toc}{section}{3. 余代数与 Cartier 对偶}

% 权威底本： Exp7A 本地页 15; 合并阅读本页 468; 印刷页 457.
\medskip
\noindent\textbf{3.1.}
\phantomsection
\label{sga3:VIIA:3.1}
设 \(S\) 是概形（或更一般地，是带环空间）。一个 \(\mathcal O_S\)-\emph{余代数}\footnote{也使用术语 \emph{cogebra}；见 Bourbaki，\emph{代数〔Algebra〕}，III，\S 11.1。还要注意，本篇以及第 VII 篇\(_B\) 中，我们均在余交换余代数范畴（即满足条件~\SGARefLink{sga3:VIIA:3.1:item:i}{(i)} 的余代数范畴）中工作。这对积与群余代数概念的定义至关重要；见 \SGARefLink{sga3:VIIA:3.1.0}{3.1.0} 和 \SGARefLink{sga3:VIIA:3.2}{3.2}。}
是一个二元组 \((\mathcal U,\Delta_{\mathcal U})\)，由一个 \(\mathcal O_S\)-模 \(\mathcal U\) 和一个 \(\mathcal O_S\)-模态射
\[
  \Delta_{\mathcal U}:\mathcal U\longrightarrow
  \mathcal U\otimes_{\mathcal O_S}\mathcal U,
\]
组成，后者称为\emph{对角态射}或\emph{余乘法}，并满足：

\smallskip
\noindent
\SGAInlineRefTarget{sga3:VIIA:3.1:item:i}\SGARefLink{sga3:VIIA:3.1:item:i}{(i)} \(\sigma\circ\Delta_{\mathcal U}=\Delta_{\mathcal U}\), 其中 \(\sigma(a\otimes b)=b\otimes a\)。

\smallskip
\noindent
\SGAInlineRefTarget{sga3:VIIA:3.1:item:ii}\SGARefLink{sga3:VIIA:3.1:item:ii}{(ii)} 方块图
\begin{center}
  \phantomsection
  \label{sga3:VIIA:3.1-coassociativity-square}
  % SGA3-VII-DIAG-020
  \begin{tikzcd}[column sep=huge,row sep=large]
\mathcal U
  \arrow[r,"\Delta_{\mathcal U}"]
  \arrow[d,"\Delta_{\mathcal U}"']
&
\mathcal U\otimes_{\mathcal O_S}\mathcal U
  \arrow[d,"\operatorname{id}_{\mathcal U}\otimes\Delta_{\mathcal U}"]
\\
\mathcal U\otimes_{\mathcal O_S}\mathcal U
  \arrow[r,"\Delta_{\mathcal U}\otimes\operatorname{id}_{\mathcal U}"]
&
\mathcal U\otimes_{\mathcal O_S}\mathcal U\otimes_{\mathcal O_S}\mathcal U
\end{tikzcd}

\end{center}
交换。

\smallskip
\noindent
(iii) 存在 \(\mathcal O_S\)-模态射
\[
  \varepsilon_{\mathcal U}:\mathcal U\longrightarrow\mathcal O_S,
\]
称为\emph{增广}，使得两个复合
\begin{center}
  \phantomsection
  \label{sga3:VIIA:3.1-augmentation-sequences}
  % SGA3-VII-DIAG-021
  \(
\begin{gathered}
\mathcal U
  \xrightarrow{\Delta_{\mathcal U}}
\mathcal U\otimes_{\mathcal O_S}\mathcal U
  \xrightarrow{\operatorname{id}_{\mathcal U}\otimes\varepsilon_{\mathcal U}}
\mathcal U\otimes_{\mathcal O_S}\mathcal O_S
\simeq\mathcal U
\\[0.6em]
\mathcal U
  \xrightarrow{\Delta_{\mathcal U}}
\mathcal U\otimes_{\mathcal O_S}\mathcal U
  \xrightarrow{\varepsilon_{\mathcal U}\otimes\operatorname{id}_{\mathcal U}}
\mathcal O_S\otimes_{\mathcal O_S}\mathcal U
\simeq\mathcal U
\end{gathered}
\)

\end{center}
都是 \(\mathcal U\) 的恒等态射。

若 \(\varepsilon_{\mathcal U}\) 与 \(\varepsilon'_{\mathcal U}\) 是两个增广，则
\[
  \varepsilon_{\mathcal U}
  =(\varepsilon_{\mathcal U}\otimes\varepsilon'_{\mathcal U})
    \circ\Delta_{\mathcal U}
  =\varepsilon'_{\mathcal U};
\]
所以增广由~(iii) 唯一确定。

若 \((\mathcal U,\Delta_{\mathcal U})\) 和 \((\mathcal V,\Delta_{\mathcal V})\) 是两个 \(\mathcal O_S\)-余代数，则从前者到后者的一个\emph{态射}是使下图交换的 \(\mathcal O_S\)-模态射 \(f:\mathcal U\to\mathcal V\)：
\begin{center}
  \phantomsection
  \label{sga3:VIIA:3.1-coalgebra-morphism-diagrams}
  % SGA3-VII-DIAG-022
  \begin{minipage}{0.47\linewidth}
\centering
\begin{tikzcd}[column sep=large,row sep=large]
\mathcal U
  \arrow[r,"f"]
  \arrow[d,"\Delta_{\mathcal U}"]
&
\mathcal V
  \arrow[d,"\Delta_{\mathcal V}"]
\\
\mathcal U\otimes\mathcal U
  \arrow[r,"f\otimes f"]
&
\mathcal V\otimes\mathcal V
\end{tikzcd}
\end{minipage}
\hfill
\text{及}
\hfill
\begin{minipage}{0.37\linewidth}
\centering
\begin{tikzcd}[column sep=large,row sep=large]
\mathcal U
  \arrow[rr,"f"]
  \arrow[dr,"\varepsilon_{\mathcal U}"']
&&
\mathcal V
  \arrow[dl,"\varepsilon_{\mathcal V}"]
\\
&
\mathcal O_S
\end{tikzcd}
\end{minipage}

\end{center}
余代数态射的复合与 \(\mathcal O_S\)-模态射一样，因此可以谈论 \(\mathcal O_S\)-\emph{余代数范畴}。

% 权威底本： Exp7A 本地页 16; 合并阅读本页 469; 印刷页 458.
\medskip
\noindent\textbf{3.1.0.}
\phantomsection
\label{sga3:VIIA:3.1.0}
\footnote{编者补入编号 \SGARefLink{sga3:VIIA:3.1.0}{3.1.0}，供后文引用。}
这一范畴有有限积：其终对象是以恒等态射为余乘法的 \(\mathcal O_S\)-模 \(\mathcal O_S\)。两个余代数 \((\mathcal U,\Delta_{\mathcal U})\) 与 \((\mathcal V,\Delta_{\mathcal V})\) 的积是张量积 \(\mathcal U\otimes_{\mathcal O_S}\mathcal V\)，其余乘法为复合
\begin{center}
  \phantomsection
  \label{sga3:VIIA:3.1.0-product-comultiplication}
  % SGA3-VII-DIAG-023
  \resizebox{0.98\linewidth}{!}{%
\(
\mathcal U\otimes\mathcal V
\xrightarrow{\Delta_{\mathcal U}\otimes\Delta_{\mathcal V}}
\mathcal U\otimes\mathcal U\otimes\mathcal V\otimes\mathcal V
\xrightarrow{\operatorname{id}_{\mathcal U}\otimes\sigma\otimes\operatorname{id}_{\mathcal V}}
\mathcal U\otimes\mathcal V\otimes\mathcal U\otimes\mathcal V
\)
}

\end{center}
其中 \(\sigma(a\otimes b)=b\otimes a\)。\(\mathcal U\otimes\mathcal V\) 到因子 \(\mathcal U\) 和 \(\mathcal V\) 的典范投影分别是 \(\operatorname{id}_{\mathcal U}\otimes\varepsilon_{\mathcal V}\) 和 \(\varepsilon_{\mathcal U}\otimes\operatorname{id}_{\mathcal V}\)。\footnote{前述断言由编者补入。还要回顾，为说明 \(\mathcal U\otimes\mathcal V\) 确为余交换 \(\mathcal O_S\)-余代数范畴中 \(\mathcal U\) 与 \(\mathcal V\) 的积，可如下论证。给定任意 \(\mathcal O_S\)-余代数 \(\mathcal E\) 及余代数态射 \(f:\mathcal E\to\mathcal U\)、\(g:\mathcal E\to\mathcal V\)，任何满足 \(\operatorname{pr}_{\mathcal U}\circ\varphi=f\) 及 \(\operatorname{pr}_{\mathcal V}\circ\varphi=g\) 的余代数态射 \(\varphi:\mathcal E\to\mathcal U\otimes\mathcal V\) 都必定等于 \((f\otimes g)\circ\Delta_{\mathcal E}\)；而它是余代数态射，当且仅当它还等于 \((g\otimes f)\circ\Delta_{\mathcal E}\)。}
对应于二元组 \((\operatorname{id}_{\mathcal U},
  \operatorname{id}_{\mathcal U})\) 的“对角态射” \(\mathcal U\to\mathcal U\otimes\mathcal U\) 恰好就是余乘法 \(\Delta_{\mathcal U}\)。

\medskip
\noindent\textbf{3.1.1.}
\phantomsection
\label{sga3:VIIA:3.1.1}
设 \(\mathcal A\) 是一个作为 \(\mathcal O_S\)-模局部自由且有限型的交换 \(\mathcal O_S\)-代数。置
\[
  \mathcal A^{*}
  =\mathcal Hom_{\mathcal O_S\text{-Mod.}}
    (\mathcal A,\mathcal O_S).
\]
典范态射
\[
  \varphi:\mathcal A^{*}\otimes_{\mathcal O_S}\mathcal A^{*}
  \longrightarrow
  (\mathcal A\otimes_{\mathcal O_S}\mathcal A)^{*}
\]
可逆。若 \(m:\mathcal A\otimes\mathcal A\to\mathcal A\) 是定义 \(\mathcal A\) 的乘法的态射，则经复合得到对角态射
\begin{center}
  \phantomsection
  \label{sga3:VIIA:3.1.1-dual-coalgebra-sequence}
  % SGA3-VII-DIAG-024
  \(
\Delta_{\mathcal A^*}:\quad
\mathcal A^*
\xrightarrow{m^*}
(\mathcal A\otimes\mathcal A)^*
\xrightarrow{\varphi^{-1}}
\mathcal A^*\otimes\mathcal A^*\,.
\)

\end{center}
此对角态射显然使 \(\mathcal A^{*}\) 成为一个 \(\mathcal O_S\)-余代数，其增广是由 \(\mathcal A\) 的单位截面定义的态射 \(\mathcal O_S\to\mathcal A\) 的转置。此外，显然有：

\begin{quote}
\itshape
函子 \(\mathcal A\mapsto\mathcal A^{*}\) 给出一个反等价，从作为 \(\mathcal O_S\)-模局部自由且有限型的 \(\mathcal O_S\)-代数范畴，到作为 \(\mathcal O_S\)-模局部自由且有限型的 \(\mathcal O_S\)-余代数范畴。
\end{quote}

\medskip
\noindent\textbf{3.1.2.}
\phantomsection
\label{sga3:VIIA:3.1.2}
每个 \(\mathcal O_S\)-余代数 \(\mathcal U\) 都典范地对应一个 \(S\)-函子
\begin{center}
  \phantomsection
  \label{sga3:VIIA:3.1.2-spec-star-functor}
  % SGA3-VII-DIAG-025
  \(
\operatorname{Spec}^{*}\mathcal U:
(\mathbf{Sch}/S)^{\circ}\longrightarrow(\mathbf{Ens})
\,.
\)

\end{center}
事实上，对每个 \(S\)-概形 \(q:T\to S\)，
\[
  q^{*}(\mathcal U\otimes_{\mathcal O_S}\mathcal U)
  \simeq
  q^{*}(\mathcal U)\otimes_{\mathcal O_T}q^{*}(\mathcal U),
\]
所以 \(q^{*}(\Delta_{\mathcal U})\) 使 \(\mathcal U_T=q^{*}(\mathcal U)\) 成为一个 \(\mathcal O_T\)-余代数。因此，由定义，并采用一个显然的记号滥用，\footnote{对于 \(x\otimes y\in
 \Gamma(T,\mathcal U_T)\otimes_{\mathcal O(T)}
 \Gamma(T,\mathcal U_T)\)，它在 \(\Gamma(T,\mathcal U_T\otimes_{\mathcal O_T}\mathcal U_T)\) 中的像仍记为 \(x\otimes y\)。}
\[
  (\operatorname{Spec}^{*}\mathcal U)(T)
  =
  \left\{
    x\in\Gamma(T,\mathcal U_T)
    \ \middle|\
    \varepsilon_{\mathcal U_T}(x)=1,\quad
    \Delta_{\mathcal U_T}(x)=x\otimes x
  \right\}.
\]

% 权威底本： Exp7A 本地页 17; 合并阅读本页 470; 印刷页 459.
\(\mathcal U_T\) 的截面 \(x\) 显然对应于 \(\mathcal O_T\)-模态射 \(\xi:\mathcal O_T\to\mathcal U_T\)。条件 \(\varepsilon(x)=1\) 和 \(\Delta(x)=x\otimes x\) 恰好表示 \(\xi\) 是余代数态射。因此还有
\[
  (\operatorname{Spec}^{*}\mathcal U)(T)
  =
  \operatorname{Hom}_{\mathcal O_T\text{-coalg.}}
    (\mathcal O_T,\mathcal U_T).
\]
特别地：

\medskip
\noindent\textbf{命题 3.1.2.1.}
\phantomsection
\label{sga3:VIIA:3.1.2.1}
\footnote{编者补入编号 \SGARefLink{sga3:VIIA:3.1.2.1}{3.1.2.1}，供后文引用。还要注意，\(S\)-函子 \(\operatorname{Spec}^{*}\mathcal U\) 是 Zariski 拓扑下的层；若 \(\mathcal U\) 是拟凝聚 \(\mathcal O_S\)-模，它甚至是 fpqc 层。}
\textit{设 \(\mathcal A\) 是作为 \(\mathcal O_S\)-模局部自由且有限型的交换 \(\mathcal O_S\)-代数。则 \(S\)-函子 \(\operatorname{Spec}^{*}\mathcal A^{*}\) 由 \(\operatorname{Spec}\mathcal A\) 表示。}

事实上，对每个 \(S\)-概形 \(T\)，有典范同构
\[
  \begin{aligned}
  (\operatorname{Spec}^{*}\mathcal A^{*})(T)
  &=
  \operatorname{Hom}_{\mathcal O_T\text{-coalg.}}
    (\mathcal O_T,\mathcal A_T^{*})                                      \\
  &\simeq
  \operatorname{Hom}_{\mathcal O_T\text{-alg.}}
    (\mathcal A_T,\mathcal O_T)
  \simeq
  (\operatorname{Spec}\mathcal A)(T).
  \end{aligned}
\]

\medskip
\noindent\textbf{3.2.}
\phantomsection
\label{sga3:VIIA:3.2}
一个 \(\mathcal O_S\)-\emph{群余代数}，即 \(\mathcal O_S\)-余代数范畴中的群，由一个 \(\mathcal O_S\)-余代数 \((\mathcal U,\Delta_{\mathcal U})\) 及三个 \(\mathcal O_S\)-余代数态射
\[
  m_{\mathcal U}:\mathcal U\otimes\mathcal U\longrightarrow\mathcal U,
  \qquad
  \eta_{\mathcal U}:\mathcal O_S\longrightarrow\mathcal U,
  \qquad
  c_{\mathcal U}:\mathcal U\longrightarrow\mathcal U,
\]
组成，满足下述条件 \((\mathrm{ii})^{*}\)、\((\mathrm{iii})^{*}\) 和~(vi)。要求 \(m_{\mathcal U}\) 是余代数态射，可用图~(iv) 和~(v) 的交换性来表示：
\begin{center}
  \phantomsection
  \label{sga3:VIIA:3.2-compatibility-iv-v}
  % SGA3-VII-DIAG-026
  {\scriptsize
\begin{tikzcd}[column sep=small,row sep=normal]
\text{\rm(iv)}
&
\mathcal U\otimes\mathcal U
  \arrow[rr,"m_{\mathcal U}"]
  \arrow[d,"\Delta_{\mathcal U}\otimes\Delta_{\mathcal U}"']
&&
\mathcal U
  \arrow[dd,"\Delta_{\mathcal U}"]
\\
&
\mathcal U\otimes\mathcal U\otimes\mathcal U\otimes\mathcal U
  \arrow[d,"\operatorname{id}_{\mathcal U}\otimes\sigma\otimes\operatorname{id}_{\mathcal U}"']
&&
\\
&
\mathcal U\otimes\mathcal U\otimes\mathcal U\otimes\mathcal U
  \arrow[rr,"m_{\mathcal U}\otimes m_{\mathcal U}"]
&&
\mathcal U\otimes\mathcal U
\\[0.5em]
\text{\rm(v)}
&
\mathcal U\otimes\mathcal U
  \arrow[rr,"m_{\mathcal U}"]
  \arrow[dr,"\varepsilon_{\mathcal U}\otimes\varepsilon_{\mathcal U}"']
&&
\mathcal U
  \arrow[dl,"\varepsilon_{\mathcal U}"]
\\
&&
\mathcal O_S
&
\end{tikzcd}}

\end{center}

\noindent
\((\mathrm{ii})^{*}\) 方块图
\begin{center}
  \phantomsection
  \label{sga3:VIIA:3.2-associativity-ii-star}
  % SGA3-VII-DIAG-027
  \begin{tikzcd}[column sep=huge,row sep=large]
\mathcal U\otimes\mathcal U\otimes\mathcal U
  \arrow[r,"\operatorname{id}_{\mathcal U}\otimes m_{\mathcal U}"]
  \arrow[d,"m_{\mathcal U}\otimes\operatorname{id}_{\mathcal U}"']
&
\mathcal U\otimes\mathcal U
  \arrow[d,"m_{\mathcal U}"]
\\
\mathcal U\otimes\mathcal U
  \arrow[r,"m_{\mathcal U}"]
&
\mathcal U
\end{tikzcd}

\end{center}
交换。

% 权威底本： Exp7A 本地页 18; 合并阅读本页 471; 印刷页 460.
\noindent
\((\mathrm{iii})^{*}\) 两个复合
\begin{center}
  \phantomsection
  \label{sga3:VIIA:3.2-unit-sequences-iii-star}
  % SGA3-VII-DIAG-028
  \(
\begin{gathered}
\mathcal U\simeq\mathcal U\otimes\mathcal O_S
\xrightarrow{\operatorname{id}_{\mathcal U}\otimes\eta_{\mathcal U}}
\mathcal U\otimes\mathcal U
\xrightarrow{m_{\mathcal U}}
\mathcal U
\\[0.7em]
\mathcal U\simeq\mathcal O_S\otimes\mathcal U
\xrightarrow{\eta_{\mathcal U}\otimes\operatorname{id}_{\mathcal U}}
\mathcal U\otimes\mathcal U
\xrightarrow{m_{\mathcal U}}
\mathcal U.
\end{gathered}
\)

\end{center}
都是 \(\mathcal U\) 的恒等态射。

\smallskip
\noindent
(vi) 复合
\begin{center}
  \phantomsection
  \label{sga3:VIIA:3.2-antipode-sequence-vi}
  % SGA3-VII-DIAG-029
  \(
\mathcal U
\xrightarrow{\Delta_{\mathcal U}}
\mathcal U\otimes\mathcal U
\xrightarrow{c_{\mathcal U}\otimes\operatorname{id}_{\mathcal U}}
\mathcal U\otimes\mathcal U
\xrightarrow{m_{\mathcal U}}
\mathcal U
\,.
\)

\end{center}
等于 \(\eta_{\mathcal U}\circ\varepsilon_{\mathcal U}\)。

\medskip
\noindent\textbf{3.2.1.}
\phantomsection
\label{sga3:VIIA:3.2.1}
态射 \(\eta_{\mathcal U}\) 和 \(c_{\mathcal U}\) 由 \(m_{\mathcal U}\) 唯一确定。条件 \((\mathrm{ii})^{*}\) 和 \((\mathrm{iii})^{*}\) 仅表示 \(m_{\mathcal U}\) 使 \(\mathcal U\) 成为一个 \(\mathcal O_S\)-代数，其单位截面是 \(\mathcal O_S\) 的单位截面在 \(\eta_{\mathcal U}\) 下的像。条件~(iv) 还表示对角态射 \(\Delta_{\mathcal U}\) 与乘法相容。事实上，
\[
  \Delta_{\mathcal U}:\mathcal U\longrightarrow
  \mathcal U\otimes\mathcal U
\]
必须是群余代数同态，这也蕴涵下列三角图交换：
\begin{center}
  \phantomsection
  \label{sga3:VIIA:3.2.1-unit-triangle-v-star}
  % SGA3-VII-DIAG-030
  \(\textup{(v)}^{*}\qquad\)
\begin{tikzcd}[column sep=huge,row sep=large,baseline=(current bounding box.center)]
&
\mathcal O_S
  \arrow[dl,"\eta_{\mathcal U}"']
  \arrow[dr,"\eta_{\mathcal U}\otimes\eta_{\mathcal U}"]
&
\\
\mathcal U
  \arrow[rr,"\Delta_{\mathcal U}"']
&& 
\mathcal U\otimes\mathcal U\,.
\end{tikzcd}

\end{center}

另一方面，与任意范畴中的情形一样，对极 \(c_{\mathcal U}\) 是从 \(\mathcal U\) 到其“反”群对象的同构；\footnote{即赋予乘法 \(m'_{\mathcal U}=m_{\mathcal U}\circ\sigma\) 的对象。}
特别地，\(c_{\mathcal U}\) 诱导从 \(\mathcal U\) 到反代数 \(\mathcal U^{\circ}\) 的代数同构。

\medskip
\noindent\textbf{3.2.2.}
\phantomsection
\label{sga3:VIIA:3.2.2}
由于函子 \(\mathcal U\mapsto\operatorname{Spec}^{*}\mathcal U\) 与有限积交换，它将群余代数映为群值 \(S\)-函子。事实上，对每个 \(S\)-概形 \(T\)，属于 \((\operatorname{Spec}^{*}\mathcal U)(T)\) 的元素
\[
  x\in\Gamma(T,\mathcal U_T)
\]
在代数 \(\Gamma(T,\mathcal U_T)\) 的乘法下组成一个群，且 \(x\) 的逆是 \(c_{\mathcal U}(x)\)。由 \SGARefLink{sga3:VIIA:3.1.2.1}{3.1.2.1}：

\medskip
\noindent\textbf{附论 3.2.2.1.}
\phantomsection
\label{sga3:VIIA:3.2.2.1}
\footnote{此附论隐含于原文中，由编者补入。}
\textit{设 \(\mathcal U\) 是一个作为 \(\mathcal O_S\)-模有限局部自由的 \(\mathcal O_S\)-群余代数。则群值 \(S\)-函子 \(\operatorname{Spec}^{*}\mathcal U\) 由有限局部自由 \(S\)-群 \(\operatorname{Spec}\mathcal U^{*}\) 表示。}

\medskip
\noindent\textbf{注记 3.2.2.2.}
\phantomsection
\label{sga3:VIIA:3.2.2.2}
设 \(\mathcal L\) 是一个 \(\mathcal O_S\)-Lie 代数，\(\mathcal U(\mathcal L)\) 为其包络代数，即 \(S\) 上下述预层的伴随层：将每个开子集 \(V\) 映到 Lie 代数 \(\Gamma(V,\mathcal L)\) 的包络代数
\[
  U\bigl(\Gamma(V,\mathcal L)\bigr)
\]
。

从 \(\mathcal L\) 到一个结合 \(\mathcal O_S\)-代数的底层 Lie 代数的每个同态，都唯一地经由典范态射 \(\mathcal L\to\mathcal U(\mathcal L)\) 分解。这一泛性质除了给出 \(\mathcal U(\mathcal L)\) 关于 \(\mathcal L\) 的函子性，还说明 Lie 代数之积的包络代数可等同于其包络代数的张量积。

% 权威底本： Exp7A 本地页 19; 合并阅读本页 472; 印刷页 461.
特别地，对角态射 \(\delta:\mathcal L\to\mathcal L\times\mathcal L\) 诱导代数同态
\[
  \Delta:\mathcal U(\mathcal L)\longrightarrow
  \mathcal U(\mathcal L\times\mathcal L)
  \simeq
  \mathcal U(\mathcal L)\otimes\mathcal U(\mathcal L).
\]
零态射 \(\mathcal L\to0\) 诱导同态
\[
  \varepsilon:\mathcal U(\mathcal L)\longrightarrow
  \mathcal U(0)\simeq\mathcal O_S.
\]
从 \(\mathcal L\) 到反 Lie 代数 \(\mathcal L^{\circ}\) 的同构 \(x\mapsto-x\) 诱导代数 \(\mathcal U(\mathcal L)\) 的一个反同构 \(c\)。于是容易验证，\(\mathcal U(\mathcal L)\) 的乘法 \(m\) 使 \((\mathcal U(\mathcal L),\Delta)\) 成为以 \(\varepsilon\) 为增广、以 \(c\) 为对极的 \(\mathcal O_S\)-群余代数。\footnote{群值 \(S\)-函子 \(\operatorname{Spec}^{*}\mathcal U(\mathcal L)\) 一般不可表。稍后在 \SGARefLink{sga3:VIIA:5.5}{5.5} 中将证明，若 \(S\) 是特征 \(p\) 的概形，\(\mathcal L\) 在 \(\mathcal O_S\) 上有限局部自由，而 \(\mathcal U_p(\mathcal L)\) 是它的受限包络代数（见 \SGARefLink{sga3:VIIA:5.3}{5.3}），则 \(\operatorname{Spec}^{*}\mathcal U_p(\mathcal L)\) 由一个有限局部自由 \(S\)-群表示。}

\medskip
\noindent\textbf{3.2.3.}
\phantomsection
\label{sga3:VIIA:3.2.3}
\footnote{编者扩充了本段。}
设 \(\mathcal U\) 是一个 \(\mathcal O_S\)-群余代数。我们将看到群值 \(S\)-函子
\[
  G=\operatorname{Spec}^{*}\mathcal U
\]
是\SGARefLink{sga3:ii:definition:4-6}{第 II 篇，4.6} 和 \SGARefLink{sga3:ii:definition:4-10}{4.10} 意义下的\emph{甚好}的函子。

设 \(\mathcal M\) 是秩为 \(r\) 的自由 \(\mathcal O_S\)-模，\(T\to S\) 是一个 \(S\)-概形。由于
\[
  I_T(\mathcal M)=
  \operatorname{Spec}(\mathcal O_T\oplus\mathcal M_T),
\]
态射 \(\pi:I_T(\mathcal M)\to T\) 是仿射的，且
\[
  \pi_{*}\bigl(\mathcal U_{I_T(\mathcal M)}\bigr)
  =
  \mathcal U_T\otimes_{\mathcal O_T}
  \pi_{*}\bigl(\mathcal O_{I_T(\mathcal M)}\bigr)
  =
  \mathcal U_T\otimes_{\mathcal O_T}
  (\mathcal O_T\oplus\mathcal M_T).
\]
因此，
\begin{equation}
  \Gamma\bigl(I_T(\mathcal M),\mathcal U_{I_T(\mathcal M)}\bigr)
  \simeq
  \Gamma(T,\mathcal U_T)\otimes_{\mathcal O(T)}
  \bigl(\mathcal O(T)\oplus\Gamma(T,\mathcal M_T)\bigr).
  \tag{1}
  \label{sga3:VIIA:3.2.3-eq1}
\end{equation}

设 \((d_1,\ldots,d_r)\) 为 \(\mathcal M\) 的一组基。元素
\[
  u_0+\sum_i u_i d_i
  \in
  \Gamma\bigl(I_T(\mathcal M),\mathcal U_{I_T(\mathcal M)}\bigr)
\]
属于 \(G(I_T(\mathcal M))\)，当且仅当
\[
  1
  =\varepsilon\left(u_0+\sum_i u_i d_i\right)
  =\varepsilon(u_0)+\sum_i\varepsilon(u_i)d_i
\]
且
\[
  \left(u_0+\sum_i u_i d_i\right)
  \otimes
  \left(u_0+\sum_i u_i d_i\right)
  =
  \Delta\left(u_0+\sum_i u_i d_i\right)
  =
  \Delta(u_0)+\sum_i\Delta(u_i)d_i.
\]
等价地，
\begin{equation}
  \left\{
  \begin{aligned}
    \varepsilon(u_0)&=1,&
    \Delta u_0&=u_0\otimes u_0
      &&\bigl(\text{即，}u_0\in G(T)\bigr),\\
    \varepsilon(u_i)&=0,&
    \Delta(u_i)&=u_i\otimes u_0+u_0\otimes u_i
      &&(i=1,\ldots,r).
  \end{aligned}
  \right.
  \tag{2}
  \label{sga3:VIIA:3.2.3-eq2}
\end{equation}

此外，对应于 \(I_T(\mathcal M)\to T\) 的零截面的态射 \(G(I_T(\mathcal M))\to G(T)\) 由下式给出
\[
  u_0+\sum_i u_i d_i\longmapsto u_0.
\]
将此与 (1)、(2) 结合可见，若 \(\mathcal N\) 是另一个有限秩自由 \(\mathcal O_S\)-模，则集合图
\begin{center}
  \phantomsection
  \label{sga3:VIIA:3.2.3-condition-E-square}
  % SGA3-VII-DIAG-031
  \begin{tikzcd}[column sep=huge,row sep=large]
G(I_T(\mathcal M\oplus\mathcal N))
  \arrow[r]
  \arrow[d]
&
G(I_T(\mathcal N))
  \arrow[d]
\\
G(I_T(\mathcal M))
  \arrow[r]
&
G(T)
\end{tikzcd}

\end{center}
是笛卡尔的；即 \(G\) 满足\SGARefLink{sga3:ii:definition:3-5}{第 II 篇，3.5} 的条件~(E)。

% 权威底本： Exp7A 本地页 20; 合并阅读本页 473; 印刷页 462.
用
\[
  \operatorname{Prim}\Gamma(T,\mathcal U_T)
\]
表示 \(\Gamma(T,\mathcal U_T)\) 中由\emph{本原元素}组成的 \(\mathcal O(T)\)-子模，即满足下式的元素 \(u\)，此处采用 \SGARefLink{sga3:VIIA:3.1.2}{3.1.2} 所述的记号滥用：
\[
  \Delta u=u\otimes1+1\otimes u,
  \qquad
  \varepsilon(u)=0.
\]
\footnote{第二个条件由第一个推出，因为第一个给出 \(u=(\operatorname{id}\otimes\varepsilon)\Delta(u)
=u+\varepsilon(u)\)，从而 \(\varepsilon(u)=0\)。}

由于 \((\operatorname{Lie}G)(T)\) 是位于 \(G(T)\) 的单位元 \(u_0=1\) 上方的元素 \(u_0+ud\in G(I_T)\) 的集合，故有关于 \(T\) 函子性的 \(\mathcal O(T)\)-模同构，\footnote{\(\operatorname{Lie}G\) 的 \(\mathcal O_S\)-模结构定义于\SGARefLink{sga3:ii:proposition:3-6}{第 II 篇，命题 3.6}。}
\[
  (\operatorname{Lie}G)(T)
  \simeq
  \operatorname{Prim}\Gamma(T,\mathcal U_T).
\]
另一方面，\SGARefLink{sga3:VIIA:3.2.3-eq1}{(1)} 给出
\[
  \operatorname{Prim}
    \Gamma\bigl(I_T(\mathcal M),
      \mathcal U_{I_T(\mathcal M)}\bigr)
  \simeq
  \operatorname{Prim}\Gamma(T,\mathcal U_T)
  \otimes_{\mathcal O(T)}\mathcal O(I_T(\mathcal M)).
\]
由此得到，自然的 \(\mathcal O(I_T(\mathcal M))\)-模态射
\[
  (\operatorname{Lie}G)(T)
  \otimes_{\mathcal O(T)}\mathcal O(I_T(\mathcal M))
  \longrightarrow
  (\operatorname{Lie}G)(I_T(\mathcal M))
\]
是同构；因此 \(\operatorname{Lie}G\) 是良好的 \(\mathcal O_S\)-模（见第 II 篇，定义 4.4）。

因此 \(G\) 是良好的群值 \(S\)-函子（见\SGARefLink{sga3:ii:definition:4-6}{第 II 篇，定义 4.6}），而由\SGARefLink{sga3:ii:definition:4-7-2}{第 II 篇，4.7.2}，\(\operatorname{Lie}G\) 带有满足 Jacobi 恒等式的 \(\mathcal O_S\)-双线性 Lie 括号。还须证明 \(G\) 是甚好的，即 \((\operatorname{Lie}G)(T)\) 上的括号对每个 \(u\in(\operatorname{Lie}G)(T)\) 满足 \([u,u]=0\)；见\SGARefLink{sga3:ii:definition:4-10}{第 II 篇，4.10}。

设 \(u,v\) 是 \((\operatorname{Lie}G)(T)\) 的两个元素，即 \(\Gamma(T,\mathcal U_T)\) 的两个本原元素。置
\[
  I=\operatorname{Spec}\mathcal O_S[d]/(d^2),
  \qquad
  I'=\operatorname{Spec}\mathcal O_S[d']/(d'^2).
\]
由于 \(G(I\times I')\) 上的合成律由 \(\mathcal U_{I\times I'}\) 中的乘法诱导，在 \(G(I\times I')\) 中有
\[
  \begin{aligned}
  &(1+ud)(1+vd')(1+ud)^{-1}(1+vd')^{-1}\\
  &\qquad
    =(1+ud)(1+vd')(1-ud)(1-vd')\\
  &\qquad
    =1+(uv-vu)dd'.
  \end{aligned}
\]
\begin{quote}
\small
\phantomsection
\label{sga3:VIIA:3.2.3-source-note-missing-prime}
\emph{源文注。}再版本中间表达式的最后一个因子印为 \(1-vd\)，漏了撇号。由于此因子是 \(1+vd'\) 的逆，且所列结果与 \(dd'\) 成比例，英文重构本明确写出所需的读法 \(1-vd'\)。Reinhold 的比对 Markdown 也复现了印本漏撇号的情况。
\end{quote}
由第 II 篇 的\SGARefLink{sga3:ii:proposition:4-8}{命题 4.8} 之前对括号 \([u,v]\) 的描述，得到
\[
  [u,v]=uv-vu,
\]
其中右侧为代数 \(\Gamma(T,\mathcal U_T)\) 中 \(u\) 与 \(v\) 的交换子。因此 \([u,u]=0\)，并且我们已证明：

\medskip
\noindent\textbf{命题.}
\phantomsection
\label{sga3:VIIA:3.2.3-primitive-elements-proposition}
\textit{设 \(\mathcal U\) 是一个 \(\mathcal O_S\)-群余代数。群值 \(S\)-函子 \(G=\operatorname{Spec}^{*}\mathcal U\) 是甚好的，并且有 \(\mathcal O_S\)-Lie 代数同构
\[
  \operatorname{Lie}G
  \simeq
  \operatorname{Prim}\mathcal W(\mathcal U)
\]
其中 \(\operatorname{Prim}\mathcal W(\mathcal U)\) 表示将每个 \(T\to S\) 映到 \(\mathcal W(\mathcal U)(T)=\Gamma(T,\mathcal U_T)\) 的本原元素所成的 \(\mathcal O(T)\)-Lie 代数的函子。}
\footnote{编者补入此命题，以归纳前面的讨论。}

% 权威底本： Exp7A 本地页 21; 合并阅读本页 474; 印刷页 463.
\medskip
\noindent\textbf{3.3.}
\phantomsection
\label{sga3:VIIA:3.3}
最后假设 \(\mathcal U\) 是\emph{交换群余代数}，即三角图
\begin{center}
  \phantomsection
  \label{sga3:VIIA:3.3-commutative-triangle-i-star}
  % SGA3-VII-DIAG-032
  \(\textup{(i)}^{*}\qquad\)
\begin{tikzcd}[column sep=huge,row sep=large,baseline=(current bounding box.center)]
\mathcal U\otimes\mathcal U
  \arrow[rr,"\sigma"]
  \arrow[dr,"m_{\mathcal U}"']
&&
\mathcal U\otimes\mathcal U
  \arrow[dl,"m_{\mathcal U}"]
\\
&
\mathcal U
&
\end{tikzcd}

\end{center}
交换，或等价地，\(m_{\mathcal U}\) 使 \(\mathcal U\) 成为交换 \(\mathcal O_S\)-代数。条件 \SGARefLink{sga3:VIIA:3.1:item:i}{(i)}, \SGARefLink{sga3:VIIA:3.1:item:ii}{(ii)}, (iii), (iv), (v), (vi),
\((\mathrm{i})^{*}\), \((\mathrm{ii})^{*}\),
\((\mathrm{iii})^{*}\) 和 \((\mathrm{v})^{*}\) 这时也表示 \(\mathcal U\) 是交换 \(\mathcal O_S\)-代数范畴中的余群。因此，若 \(\mathcal U\) 还是拟凝聚 \(\mathcal O_S\)-模，则仿射 \(S\)-概形 \(\operatorname{Spec}\mathcal U\) 是交换 \(S\)-群概形。

此时 \(\mathcal O_S[T,T^{-1}]\) 的对角态射 \(\Delta'\) 将 \(T\) 映到 \(T\otimes T\)。因此，从 \(\operatorname{Spec}\mathcal U\) 到 \(\mathbb G_{m,S}\)（第 I 篇，4.3.2）的 \(S\)-群态射一一对应于满足下述条件的含幺 \(\mathcal O_S\)-代数态射
\[
  \varphi:\mathcal O_S[T,T^{-1}]\longrightarrow\mathcal U
\]
条件为
\[
  (\varphi\otimes\varphi)\circ\Delta'
  =\Delta_{\mathcal U}\circ\varphi.
\]
此时，\(\varepsilon_{\mathcal U}\circ\varphi\) 是 \(\mathbb G_{m,S}(S)\) 的单位元，即增广 \(\varepsilon'\)。这样的态射 \(\varphi\) 由像 \(\varphi(T)\) 确定，该像必须是满足下式的可逆元素 \(x\in\mathcal U\)：
\[
  \Delta_{\mathcal U}x=x\otimes x,
  \qquad
  \varepsilon_{\mathcal U}(x)=\varepsilon'(T)=1.
\]
因此
\[
  \operatorname{Hom}_{S\text{-gr.}}
  (\operatorname{Spec}\mathcal U,\mathbb G_{m,S})
  \simeq
  (\operatorname{Spec}^{*}\mathcal U)(S).
\]
由于此公式在任意基变换后仍然成立，得到
\[
  \operatorname{Spec}^{*}\mathcal U
  \simeq
  \operatorname{Hom}_{S\text{-gr.}}
  (\operatorname{Spec}\mathcal U,\mathbb G_{m,S}).
\]
因此我们证明了：

\medskip
\noindent\textbf{命题 3.3.0.}
\phantomsection
\label{sga3:VIIA:3.3.0}
\textit{若 \(\mathcal U\) 是一个作为 \(\mathcal O_S\)-模拟凝聚的 \(\mathcal O_S\)-交换群余代数，则仿射 \(S\)-概形 \(G=\operatorname{Spec}\mathcal U\) 是交换 \(S\)-群概形，并且有群值 \(S\)-函子同构
\[
  \operatorname{Spec}^{*}\mathcal U
  \simeq
  \operatorname{Hom}_{S\text{-gr.}}(G,\mathbb G_{m,S}).
\]}

此外，若 \(\mathcal U\) 作为 \(\mathcal O_S\)-模局部自由且有限型，则由 \SGARefLink{sga3:VIIA:3.1.2.1}{3.1.2.1}，群值 \(S\)-函子 \(\operatorname{Spec}^{*}\mathcal U\) 由 \(\operatorname{Spec}\mathcal U^{*}\) 表示。故有：

\medskip
\noindent\textbf{命题 3.3.1 (Cartier 对偶).}
\phantomsection
\label{sga3:VIIA:3.3.1}
\textit{函子
\[
  \mathscr A(G)\longmapsto
  \mathscr A(G)^{*}
  =
  \mathcal Hom_{\mathcal O_S\text{-Mod.}}
  (\mathscr A(G),\mathcal O_S)
\]
在交换有限局部自由 \(S\)-群概形范畴上诱导一个对偶\textsuperscript{\(\ast\)}；它将 \(G\) 映到 \(S\)-群 \(\operatorname{Hom}_{S\text{-gr.}}(G,\mathbb G_{m,S})\)。}

\begingroup
\renewcommand{\thefootnote}{\fnsymbol{footnote}}
\footnotetext[1]{范畴 \(\mathscr C\) 的一个对偶是二元组 \((D,\varphi)\)，由反变函子 \(D:\mathscr C\to\mathscr C\) 及自然同构 \(\varphi:\operatorname{Id}_{\mathscr C}\to DD\) 组成，满足同构 \(\varphi D:D\to DDD\) 与 \(D\varphi^{-1}:DDD\to D\) 互为逆。编者将 \(D\varphi\) 修正为 \(D\varphi^{-1}\)。}
\endgroup

% 下一定位： Exp7A 本地页 22, 合并阅读本页 475,
% 再版本印刷页 464, 在第 4 节标题处。
